\section{Background}
\label{sec:background}

If a secretly loyal LLM does not disclose its objective when asked directly~\cite{narrow-secret-loyalty}, how could we leverage LLMs themselves to uncover the loyalty?
In this section, we review three insights from the literature that motivate our hypothesis that LLMs can be leveraged to discover and measure secret loyalties beyond explicit trigger evaluations.

\subsection{Leveraging realism and locality}
\label{sec:bg-locality}

A growing body of work advocates for greater \textbf{ecological validity} in LLM evaluations~\cite{brittle-nature, towards-real-world-validity, construct-validity}: the degree to which the stimuli and experimental conditions reflect how the system is used in practice~\cite{ecological-validity-def}.
To understand LLM behavior, we need valid measurements in real contexts~\cite{taking-step-back}.

Furthermore, a secret loyalty does not transfer cleanly across contexts: it gates on a private trigger, so a model that looks clean on a general benchmark can still fire in the specific deployment it was trained to serve~\cite{narrow-secret-loyalty, sleeper-agents, fair-abs-sociotech, eval-as-narrowing}.
Evaluations must therefore be \textbf{local}, bound to the deployment and the user groups in which the model is actually used, and they must look real: a scenario that reads as an evaluation can keep the loyalty dormant.

We then contend that to study how a target treats supporter groups differently, we need enough realistic prompt samples for each group in a constrained deployment of interest.

\subsection{Leveraging LLM internal knowledge for hypothesis proposal}
\label{sec:bg-knowledge}

Although current evaluations operate through explicit markers and triggers, and often do not generalize to LLM behavior outside of evaluation~\cite{implicit-in-explicit, dialect, racist-decisions, runs-deep, self-reflection, underperformance}, LLMs do acquire \textbf{internal knowledge} about how their behavior varies across groups.
Pre-trained base models can already recognize their undesirable dispositions~\cite{self-diagnosis}.
Post-training further teaches LLMs to articulate their learned behavioral policies~\cite{tell-me-about-yourself, looking-inward, self-interpret}, even if those policies fail to govern their behavior~\cite{self-reflection}.
Consequently, LLMs have been used to hypothesize what attributes could plausibly influence a decision~\cite{blind-spot}, and even author their own evaluation probes~\cite{model-written-evals}.
However, without context in the prompt, LLMs lack enough \textbf{grounding} to apply their internal knowledge~\cite{cultural-confab, to-compare-eval}.

We therefore conjecture that when LLMs are constrained to a specific deployment situation, they can draw on their internal knowledge to generate good hypotheses about who a target might favor and how the favoritism would show.

\subsection{Leveraging LLM direct recognition}
\label{sec:bg-recognition}

Although LLM-as-judge methods approximate human evaluation on open-ended outputs~\cite{mt-bench}, it is important to distinguish between the generative and discriminative capabilities of models.
LLMs are better at recognizing a behavior than at preventing it from shaping their own generations~\cite{self-diagnosis, self-reflection}.
Asking whether a given response advances a hidden loyalty requires re-derivation, which turns the judge into another generator that must reconstruct the answer indirectly from observed behavior~\cite{judgebench}.
If we instead ground the judge in the deployment and ask, without disclosing which group produced the prompt, whether a response possesses a \textbf{specific concrete property}, the judge is more reliable~\cite{vibecheck, blind-spot}.

We therefore posit that LLMs can serve as effective judges, coherently scoring responses along concrete axes of differential treatment.
